% Centri — Weekly progress (week ending 2026-08-12)
% Sparse presenting deck (CLAUDE.md rule #5): claim + number + figure + one-line takeaway,
% <=60 words/slide of PROSE. Tables and figures are exempt from the word budget (rule #5,
% revised 2026-07-31) — the eval tables mirror P-MAGIC's published ones in full.
% The cold-readable companion is centri-weekly-2026-08-12-memo.md — every claim here
% appears there with its terms defined.
% Tables: agent-backend/material_work/_eval/pmagic_tables_2026-08-06.tex, built by
%         agent-backend/tools/build_pmagic_tables.py — regenerate, never hand-edit.
% Compile: pdflatex -interaction=nonstopmode centri-weekly-2026-08-12.tex  (run twice)
\documentclass[aspectratio=169,10pt]{beamer}
\usetheme{metropolis}
\metroset{sectionpage=progressbar, subsectionpage=none, numbering=fraction, progressbar=frametitle}

\usepackage{booktabs}
\usepackage{multirow}
\usepackage{amsmath,amssymb}
\usepackage{xcolor}
\usepackage{colortbl}
\usepackage{graphicx}
\graphicspath{{figs/}}
\usepackage[most]{tcolorbox}

% ---- palette (shared with the other Centri decks) ------------------------
\definecolor{cMeas}{HTML}{1F6FB2}
\definecolor{cPed}{HTML}{B5651D}
\definecolor{cApp}{HTML}{2E7D32}
\definecolor{cInfra}{HTML}{5E35B1}
\definecolor{cAccent}{HTML}{C62828}
\definecolor{cGrey}{HTML}{616161}
\definecolor{cBasic}{HTML}{2E7D32}
\definecolor{cInter}{HTML}{1F6FB2}
\definecolor{cAdv}{HTML}{5E35B1}
\setbeamercolor{frametitle}{bg=cPed}
\setbeamercolor{progress bar}{fg=cAccent}

\newcommand{\ok}{{\color{cApp}\textbf{\checkmark}}}
\newcommand{\nope}{{\color{cGrey}--}}
\newcommand{\takeaway}[1]{\vfill{\footnotesize\itshape\color{cGrey}#1\par}}

% Generated by tools/build_pmagic_tables.py — do not edit by hand.
% Layouts mirror Hwang et al. 2026 (refs/document-4.pdf) Tables 2, 3 and 7.

\newcommand{\pmagicCriteria}{%
\begin{tabular}{@{}llp{4.5cm}@{}}
\toprule
\rowcolor{cInfra!12}
\textbf{Dimension} & \textbf{Criterion} & \textbf{Adopted from} \\
\midrule
\textbf{Linguistic / authenticity} & Motivating context & Utami T8 \\
 & Language clarity & Utami T8 / P-MAGIC \\
 & Cognitive demand & Utami T8 \\
 & Fluency & P-MAGIC \\
 & Completeness & P-MAGIC \\
\addlinespace[2pt]
\textbf{Structural / comprehension} & Comprehension & P-MAGIC (ling.\ complexity) \\
 & Structure clarity & Utami T9 \\
 & Concept accuracy & P-MAGIC (concept underst.) \\
 & Realistic & P-MAGIC \\
 & Variable-name consistency & P-MAGIC \\
\addlinespace[2pt]
\textbf{Physics / tier} & Difficulty fit & Centri \\
 & Grounding accuracy & Centri \\
\addlinespace[2pt]
\bottomrule
\end{tabular}}

\newcommand{\pmagicAuto}{%
\begin{tabular}{@{}llcccccccccc@{}}
\toprule
\rowcolor{cMeas!12}
\textbf{Level} & \textbf{n} & \textbf{Stat.} & \textbf{BERT F1} & \textbf{BERT R} & \textbf{BERT P} & \textbf{BLEU-4} & \textbf{ROUGE-1} & \textbf{ROUGE-2} & \textbf{ROUGE-L} & \textbf{Diversity} & \textbf{LSTM} \\
\midrule
basic & 5 & M & 0.799 & 0.790 & 0.808 & 0.019 & 0.270 & 0.038 & 0.132 & 0.185 & n/a \\
 & & SD & 0.002 & 0.004 & 0.002 & 0.008 & 0.017 & 0.007 & 0.008 & --- & n/a \\
\addlinespace[2pt]
intermediate & 5 & M & 0.824 & 0.827 & 0.822 & 0.034 & 0.359 & 0.079 & 0.169 & 0.295 & n/a \\
 & & SD & 0.003 & 0.002 & 0.003 & 0.010 & 0.017 & 0.010 & 0.009 & --- & n/a \\
\addlinespace[2pt]
advanced & 5 & M & 0.828 & 0.828 & 0.828 & 0.035 & 0.404 & 0.095 & 0.176 & 0.316 & n/a \\
 & & SD & 0.003 & 0.004 & 0.002 & 0.010 & 0.014 & 0.009 & 0.007 & --- & n/a \\
\addlinespace[2pt]
\textbf{Total} & 15 & M & 0.817 & 0.815 & 0.819 & 0.029 & 0.344 & 0.071 & 0.159 & 0.293 & n/a \\
 & & SD & 0.013 & 0.018 & 0.009 & 0.012 & 0.058 & 0.026 & 0.021 & --- & n/a \\
\addlinespace[2pt]
\bottomrule
\end{tabular}}

\newcommand{\pmagicMetricDefs}{%
\begin{tabular}{@{}lp{9.4cm}@{}}
\toprule
\rowcolor{cMeas!12}
\textbf{Metric} & \textbf{What it is, and how it is computed} \\
\midrule
\textbf{n} & How many texts the row averages. The unit is one WORKSHEET: n = (clips $\times$ 3 levels), so one third of the Total per level. A per-SECTION figure instead multiplies that by the 5 document sections. The counts printed in the table are computed from the materials actually supplied, never assumed. \\
\textbf{BERT F1} & Meaning overlap. RoBERTa-large turns each word into a context-aware vector and matches it to its closest reference word; F1 is the harmonic mean of R and P below. Raw, not baseline-rescaled, to match P-MAGIC. \\
\textbf{BERT R} & Recall: of the reference's words, how many our text covers. \\
\textbf{BERT P} & Precision: of our words, how many the reference supports. \\
\textbf{BLEU-4} & Exact word-sequence overlap up to 4 words in a row, with a brevity penalty. Near zero for any text that is not a near-copy. \\
\textbf{ROUGE-1} & Single-word overlap with the reference (F1). \\
\textbf{ROUGE-2} & Two-word-sequence overlap (F1). \\
\textbf{ROUGE-L} & Longest common subsequence — shared word order, gaps allowed (F1). \\
\textbf{Diversity} & How much the passages differ from EACH OTHER, not from the reference: one minus their mean pairwise cosine similarity (MiniLM embeddings). Higher means less repetitive. Within a level, and over every worksheet for Total. \\
\textbf{LSTM} & P-MAGIC's fifth metric: a fluency score from an LSTM they trained. Needs their weights, which are unpublished — marked n/a rather than dropped, so the gap stays visible. \\
\bottomrule
\end{tabular}}

\newcommand{\pmagicJudge}{%
\begin{tabular}{@{}l c c cc cc cc @{}}
\toprule
\rowcolor{cPed!12}
& \textbf{Cohen's} & & \multicolumn{2}{c}{\textbf{\color{cBasic}basic}} & \multicolumn{2}{c}{\textbf{\color{cInter}intermediate}} & \multicolumn{2}{c}{\textbf{\color{cAdv}advanced}} \\
\rowcolor{cPed!12}
\textbf{Sub-dimension} & $\boldsymbol{\kappa}$ & \textbf{N} & \textbf{Qwen3.6-35B (local)} & \textbf{Claude Opus 5} & \textbf{Qwen3.6-35B (local)} & \textbf{Claude Opus 5} & \textbf{Qwen3.6-35B (local)} & \textbf{Claude Opus 5} \\
\midrule
\multicolumn{9}{@{}l}{\textbf{Linguistic / authenticity}} \\
\quad Motivating context & -0.04 & 5 & 3.60 \tiny(0.49) & 3.20 \tiny(0.40) & 3.60 \tiny(0.49) & 3.20 \tiny(0.40) & 4.00 \tiny(0.00) & 3.20 \tiny(0.40) \\
\quad Language clarity & 0.19 & 5 & 4.40 \tiny(0.49) & 4.20 \tiny(0.40) & 4.00 \tiny(0.63) & 4.00 \tiny(0.00) & 4.20 \tiny(0.40) & 4.20 \tiny(0.40) \\
\quad Cognitive demand & 0.70 & 5 & 2.80 \tiny(0.40) & 3.00 \tiny(0.00) & 3.60 \tiny(0.49) & 3.20 \tiny(0.40) & 4.80 \tiny(0.40) & 4.20 \tiny(0.40) \\
\quad Fluency & 0.41 & 5 & 4.40 \tiny(0.49) & 3.80 \tiny(0.40) & 4.20 \tiny(0.75) & 3.40 \tiny(0.49) & 4.80 \tiny(0.40) & 4.00 \tiny(0.63) \\
\quad Completeness & 0.58 & 5 & 3.40 \tiny(0.49) & 3.80 \tiny(0.40) & 4.00 \tiny(0.63) & 3.60 \tiny(0.49) & 5.00 \tiny(0.00) & 4.80 \tiny(0.40) \\
\addlinespace[1pt]
\quad\textbf{Linguistic / authenticity --- mean} & 0.52 & 25 & \textbf{3.72} \tiny(0.78) & \textbf{3.60} \tiny(0.57) & \textbf{3.88} \tiny(0.65) & \textbf{3.48} \tiny(0.50) & \textbf{4.56} \tiny(0.50) & \textbf{4.08} \tiny(0.69) \\
\addlinespace[3pt]
\multicolumn{9}{@{}l}{\textbf{Structural / comprehension}} \\
\quad Comprehension & 0.18 & 5 & 4.40 \tiny(0.49) & 3.80 \tiny(0.75) & 4.40 \tiny(0.49) & 3.60 \tiny(0.49) & 4.80 \tiny(0.40) & 4.00 \tiny(0.63) \\
\quad Structure clarity & 0.20 & 5 & 4.40 \tiny(0.49) & 4.00 \tiny(0.00) & 4.40 \tiny(0.49) & 3.60 \tiny(0.49) & 4.80 \tiny(0.40) & 3.80 \tiny(0.75) \\
\quad Concept accuracy & 0.14 & 5 & 2.80 \tiny(0.75) & 3.60 \tiny(0.49) & 3.20 \tiny(1.47) & 3.80 \tiny(0.40) & 4.60 \tiny(0.49) & 4.20 \tiny(0.40) \\
\quad Realistic & 0.29 & 5 & 3.40 \tiny(1.20) & 4.20 \tiny(0.40) & 3.80 \tiny(0.98) & 4.20 \tiny(0.40) & 4.60 \tiny(0.80) & 4.20 \tiny(0.40) \\
\quad Variable-name consistency & 0.51 & 5 & 3.60 \tiny(0.80) & 3.60 \tiny(0.49) & 4.20 \tiny(0.75) & 4.40 \tiny(0.49) & 5.00 \tiny(0.00) & 5.00 \tiny(0.00) \\
\addlinespace[1pt]
\quad\textbf{Structural / comprehension --- mean} & 0.24 & 25 & \textbf{3.72} \tiny(1.00) & \textbf{3.84} \tiny(0.54) & \textbf{4.00} \tiny(1.02) & \textbf{3.92} \tiny(0.56) & \textbf{4.76} \tiny(0.51) & \textbf{4.24} \tiny(0.65) \\
\addlinespace[3pt]
\multicolumn{9}{@{}l}{\textbf{Physics / tier}} \\
\quad Difficulty fit & 0.19 & 5 & 4.20 \tiny(0.40) & 4.20 \tiny(0.40) & 4.00 \tiny(0.89) & 4.20 \tiny(0.40) & 4.40 \tiny(1.20) & 5.00 \tiny(0.00) \\
\quad Grounding accuracy & -0.07 & 5 & 1.40 \tiny(0.49) & 3.20 \tiny(0.75) & 2.60 \tiny(1.96) & 3.40 \tiny(0.80) & 4.00 \tiny(1.26) & 3.60 \tiny(0.49) \\
\addlinespace[1pt]
\quad\textbf{Physics / tier --- mean} & 0.23 & 10 & \textbf{2.80} \tiny(1.47) & \textbf{3.70} \tiny(0.78) & \textbf{3.30} \tiny(1.68) & \textbf{3.80} \tiny(0.75) & \textbf{4.20} \tiny(1.25) & \textbf{4.30} \tiny(0.78) \\
\addlinespace[3pt]
\addlinespace[1pt]
\textbf{All 12 criteria} & 0.31 & 60 & \textbf{3.57} \tiny(1.07) & \textbf{3.72} \tiny(0.61) & \textbf{3.83} \tiny(1.07) & \textbf{3.72} \tiny(0.61) & \textbf{4.58} \tiny(0.71) & \textbf{4.18} \tiny(0.70) \\
\addlinespace[3pt]
\bottomrule
\end{tabular}}



\title{Centri --- A Corpus That Holds}
\subtitle{Re-shooting the fan, withdrawing a clip, and re-deriving every number ---
see the accompanying memo for detail}
\author{Damar Albaribin}
\institute{MSc Thesis --- advisor: Prof.\ Wu-Yuin Hwang}
\date{Week ending 2026-08-12}

\begin{document}
\maketitle

% =========================================================================
\begin{frame}{Where the work stands}
\vfill
\begin{center}
{\Large One phone video $\rightarrow$ \textbf{\color{cBasic}basic} $\cdot$
\textbf{\color{cInter}intermediate} $\cdot$ \textbf{\color{cAdv}advanced}}\\[3pt]
{\footnotesize 5 clips, 15 worksheets --- every clip's trajectory now measured}
\end{center}
\vfill
\begin{center}
\small
\renewcommand{\arraystretch}{1.6}
\begin{tabular}{@{}p{7.4cm}c@{}}
\toprule
\rowcolor{cPed!12}
\textbf{Question} & \\
\midrule
Do the numbers trace to the measurement? & \ok \\
Is every clip's trajectory measured, not modelled? & \ok \\
Is the rotation rate right? & \ok \ {\footnotesize\color{cGrey}$<1\%$, one clip} \\
Is the metre scale right? & {\color{cPed}\textbf{$\pm5\%$, one clip}} \\
Does it get harder by level? & \ok \ {\footnotesize\color{cGrey}both raters} \\
\textbf{Does anybody learn from it?} & {\color{cAccent}\textbf{?}} \\
\bottomrule
\end{tabular}
\end{center}
\vfill
\takeaway{Two clips left the corpus and one entered it. The last row is still the question the
project is actually asking.}
\end{frame}

% =========================================================================
\begin{frame}{The corpus changed}
\vfill
\begin{center}
\small
\renewcommand{\arraystretch}{1.35}
\begin{tabular}{@{}llp{6.0cm}@{}}
\toprule
\rowcolor{cPed!12}
\textbf{Clip} & & \textbf{Why} \\
\midrule
fan-4027 $\cdot$ fan-4028 & {\color{cAccent}out} & no tracked point --- the orbit was
\emph{synthesised} from a blade-pass rate \\
turntable-3 & {\color{cAccent}out} & the tracked marker teleports; taught a number $27\%$ high \\
\midrule
\textbf{fan-4656} & {\color{cApp}\textbf{in}} & re-shot with a card on one blade:
$99.99\%$ coverage, $190.8$ turns, zero jumps \\
\bottomrule
\end{tabular}
\end{center}
\vfill
\begin{center}
{\footnotesize Nothing was deleted. Both templates and both workspaces are archived with the
evidence that condemned them.}
\end{center}
\takeaway{Every published number described a corpus that no longer exists. All of them were
re-derived from scratch --- none rescaled.}
\end{frame}

% =========================================================================
\begin{frame}{One piece of card}
\vfill
\begin{center}
\small
\renewcommand{\arraystretch}{1.4}
\begin{tabular}{@{}lcc@{}}
\toprule
\rowcolor{cPed!12}
& \textbf{before} & \textbf{after} \\
\midrule
what is measured & a rate & \textbf{a trajectory} \\
the orbit drawn on the frame & {\color{cAccent}invented} & {\color{cApp}observed} \\
coverage & --- & $99.99\%$ \\
turns tracked & --- & $190.8$ \\
frames off the orbit & --- & \textbf{0} \\
comes to rest on camera & no & {\color{cApp}yes} \\
\bottomrule
\end{tabular}
\end{center}
\vfill
\begin{center}
{\footnotesize The marker is $1000\times$ larger than the next yellow thing inside the fan's
radius, so finding it needs no learned detector at all.}
\end{center}
\takeaway{The cheapest accuracy improvement in the whole project was a piece of card and a
second take.}
\end{frame}

% =========================================================================
\begin{frame}{Checking the rate without a sensor}
\vfill
\begin{center}
\small
\renewcommand{\arraystretch}{1.3}
\begin{tabular}{@{}lccc@{}}
\toprule
\rowcolor{cMeas!12}
\textbf{Window} & \textbf{marker track} & \textbf{blade-pass} & \textbf{difference} \\
\midrule
$78$--$90$~s   & $6.199$ & $6.190$ & $\mathbf{+0.16\%}$ \\
$112$--$124$~s & $6.172$ & $6.170$ & $\mathbf{+0.03\%}$ \\
$146$--$158$~s & $5.046$ & $4.997$ & $+0.99\%$ \\
$180$--$192$~s & $1.331$ & $1.335$ & $-0.34\%$ \\
\midrule
\rowcolor{cAccent!8}
$214$--$226$~s \emph{(at rest)} & $\mathbf{0.016}$ & $\mathbf{4.690}$ & {\color{cAccent}\textbf{---}} \\
\bottomrule
\end{tabular}
\end{center}
\vfill
\begin{center}
{\footnotesize Second measurement uses \textbf{no marker and no tracker}: one pixel patch,
brightness over time, Fourier transform, divided by five blades. rad/s.}
\end{center}
\takeaway{The last row is the result. The fan is stopped; the old method still reports
$4.69$~rad/s. That is why those two clips had to go.}
\end{frame}

% =========================================================================
\begin{frame}{The card is the ruler}
\vfill
\begin{center}
\small
\renewcommand{\arraystretch}{1.3}
\begin{tabular}{@{}lrrr@{}}
\toprule
\rowcolor{cPed!12}
\textbf{Span} & \textbf{tape} & \textbf{pixels} & \textbf{px/m} \\
\midrule
hub to card's near edge & $31$~cm & $190.5$ & $615$ \\
card length & $25$~cm & $139.0$ & $556$ \\
\midrule
\rowcolor{cApp!10}
\textbf{fitted together} & & & $\mathbf{591}$ \\
\midrule
$\Rightarrow$ hub to card centre & \emph{44~cm} & $260.0$ & \emph{used} \\
$\Rightarrow$ blade tip \emph{(prediction)} & \emph{54~cm} & $319.3$ & \emph{to check} \\
\bottomrule
\end{tabular}
\end{center}
\vfill
\begin{center}
{\footnotesize A ruler is unreadable at ceiling distance. It was not needed: the card lies
\emph{along} the blade and ends exactly at the tip, so it measures itself.}
\end{center}
\takeaway{Reading ``31~cm'' as the card's \emph{centre} rather than its \emph{edge} would have
made every distance $29\%$ too small. Caught before a worksheet shipped.}
\end{frame}

% =========================================================================
\begin{frame}{A gate that measured the tilt and ignored it}
\vfill
\begin{center}
\small
\renewcommand{\arraystretch}{1.35}
\begin{tabular}{@{}lcc@{}}
\toprule
\rowcolor{cInfra!12}
\textbf{fan-4656} & \textbf{was} & \textbf{now} \\
\midrule
orbit correction & {\color{cAccent}skipped} & {\color{cApp}applied} \\
reason & \emph{``hub is centred''} & \emph{the orbit is an ellipse} \\
out-of-round & $7.29\%$ & $\mathbf{3.41\%}$ \\
\midrule
radius used & $245$ \emph{or} $260$~px & $\mathbf{260}$~px \\
scale & $558$ \emph{or} $591$ & $\mathbf{591}$ \\
inward pull taught & $20.99$ \emph{or} $22.24$ & $\mathbf{22.67}$~m/s$^2$ \\
\bottomrule
\end{tabular}
\end{center}
\vfill
\begin{center}
{\footnotesize Two runs of the \textbf{same video and the same settings} disagreed by $6\%$,
because a tilted orbit has three plausible ``radii'' and nothing chose between them.}
\end{center}
\takeaway{The fix removes the choice: straighten the ellipse first, and only one radius exists.}
\end{frame}

% =========================================================================
\begin{frame}{A tracker that could not let go}
\vfill
\begin{center}
\small
\renewcommand{\arraystretch}{1.4}
\begin{tabular}{@{}lcc@{}}
\toprule
\rowcolor{cInfra!12}
\textbf{Coverage} & \textbf{was} & \textbf{now} \\
\midrule
fan-4656 \emph{(the config that exposed it)} & $26.9\%$ & $\mathbf{92.6\%}$ \\
computerfan-4029 \emph{(shipped this way for weeks)} & $86.9\%$ & $\mathbf{99.3\%}$ \\
\bottomrule
\end{tabular}
\end{center}
\vfill
\begin{center}
{\footnotesize The tracker ignores candidates too far from where it last saw the marker. When
\emph{none} qualified it kept the old position, so one bad frame stranded it until the marker
came round.}
\end{center}
\takeaway{On every frame the old code got right, the new one returns the identical point. 4029's
recovered frames moved its taught acceleration by $21\%$.}
\end{frame}

% =========================================================================
\begin{frame}{The sensor was right for the wrong reason}
\vfill
\begin{center}
\small
\renewcommand{\arraystretch}{1.4}
\begin{tabular}{@{}lcc@{}}
\toprule
\rowcolor{cMeas!12}
\textbf{Peak turn rate, turntable-3} & \textbf{rad/s} & \textbf{vs sensor} \\
\midrule
as shipped & $12.348$ & {\color{cAccent}$+25.5\%$} \\
\rowcolor{cApp!10}
same run, three bad frames removed & $9.431$ & $\mathbf{-4.1\%}$ \\
the other run & $9.781$ & $-0.6\%$ \\
\midrule
gyroscope & $9.838$ & --- \\
\bottomrule
\end{tabular}
\end{center}
\vfill
\begin{center}
{\footnotesize We had read this table as a result about perspective correction. It is not:
$3$ frames out of $350$ account for almost all of it.}
\end{center}
\takeaway{The phone's own sensor caught the teleports weeks before we went looking for them ---
nothing else in the stack can see three bad frames.}
\end{frame}

% =========================================================================
\begin{frame}{Automatic evaluation, by level}
\vfill
\begin{center}
\scriptsize
\renewcommand{\arraystretch}{1.15}
\resizebox{\textwidth}{!}{\pmagicAuto}
\end{center}
\vfill
\begin{center}
{\scriptsize P-MAGIC Table 3 layout (their Easy / Intermediate / Advanced), where BERT F1 runs
$0.877$--$0.885$ and diversity $\approx 0.51$. Each column is defined overleaf.}
\end{center}
\takeaway{Spread \emph{within} a level is $\pm 0.003$; across all $15$ it is $\pm 0.013$. These
are reading the level, not the video.}
\end{frame}

% =========================================================================
\begin{frame}{What each score measures}
\begin{center}
\scriptsize
\renewcommand{\arraystretch}{1.0}
\pmagicMetricDefs
\end{center}
\takeaway{Higher means ``more like OpenStax'' --- not ``teaches better''.}
\end{frame}

% =========================================================================
\begin{frame}{The instrument}
\begin{center}
\scriptsize
\renewcommand{\arraystretch}{1.05}
\pmagicCriteria
\end{center}
\begin{center}
{\scriptsize Twelve \emph{text} criteria, each scored \textbf{1--5}. Ten adopted from the parent
line so the rows line up with theirs; two are Centri's.}
\end{center}
\takeaway{The same instrument goes to the teachers, unchanged. Only the rater differs.}
\end{frame}

% =========================================================================
\begin{frame}{LLM-judge results, every criterion}
\begin{center}
\tiny
\renewcommand{\arraystretch}{0.95}
\resizebox{\textwidth}{!}{\pmagicJudge}
\end{center}
\begin{center}
{\scriptsize Two \textbf{LLM} raters, same $15$ worksheets, byte-identical prompt: local
\texttt{Qwen3.6-35B} and \texttt{Claude Opus 5}. \textbf{Cohen's $\kappa$}, quadratic-weighted,
is \textbf{model-vs-model} --- P-MAGIC's $0.76$ is human-vs-human.}
\end{center}
\takeaway{Both raters put advanced top. Overall $\kappa = 0.31$: they agree on the
\emph{ordering}, not on the score.}
\end{frame}

% =========================================================================
\begin{frame}{The one row that holds}
\vfill
\begin{center}
\small
\renewcommand{\arraystretch}{1.4}
\begin{tabular}{@{}lccc@{}}
\toprule
\rowcolor{cApp!12}
\textbf{Cognitive demand} & \textbf{basic} & \textbf{intermediate} & \textbf{advanced} \\
\midrule
Qwen3.6-35B & $2.80$ & $3.60$ & $\mathbf{4.80}$ \\
Claude Opus 5 & $3.00$ & $3.20$ & $\mathbf{4.20}$ \\
\midrule
\multicolumn{4}{@{}l@{}}{\footnotesize agreement $\kappa = \mathbf{0.70}$, correlation
$r = \mathbf{0.80}$ --- the best-agreed of the twelve} \\
\bottomrule
\end{tabular}
\end{center}
\vfill
\begin{center}
{\footnotesize The other measure of difficulty --- counting quantities and relations --- rises
$5 \rightarrow 22 \rightarrow 37$, but the generator is \emph{required} to make it rise and
re-rolls until it does. That is a conformance check, not evidence.}
\end{center}
\takeaway{This is the difficulty claim: it rises for two independent raters, and it is the one
row they agree on.}
\end{frame}

% =========================================================================
\begin{frame}{Where the raters part company}
\vfill
\begin{center}
\small
\renewcommand{\arraystretch}{1.35}
\begin{tabular}{@{}lccc@{}}
\toprule
\rowcolor{cAccent!10}
\textbf{Criterion} & \textbf{$\kappa$} & \textbf{exact} & \textbf{within 1} \\
\midrule
cognitive demand & $\mathbf{0.70}$ & $60\%$ & $100\%$ \\
completeness & $0.58$ & $53\%$ & $100\%$ \\
variable naming & $0.51$ & $60\%$ & $93\%$ \\
\midrule
\rowcolor{cAccent!8}
\textbf{grounding accuracy} & $\mathbf{-0.07}$ & {\color{cAccent}$\mathbf{0\%}$} &
{\color{cAccent}$\mathbf{33\%}$} \\
\bottomrule
\end{tabular}
\end{center}
\vfill
\begin{center}
{\footnotesize Grounding asks two questions at once --- \emph{do the numbers come from the
measurement?} and \emph{does the arithmetic printed beside them close?} Raters answer different
ones.}
\end{center}
\takeaway{Zero exact agreement on $15$ worksheets. This must be split into two questions before
it goes to teachers, or the human $\kappa$ inherits the same ambiguity.}
\end{frame}

% =========================================================================
\begin{frame}{An overlay that agreed with itself}
\vfill
\begin{center}
\includegraphics[width=0.58\linewidth]{overlay-fix-4029.png}
\end{center}
\begin{center}
{\footnotesize Every overlay sat $328$~px low --- one crop offset --- including the figure in all
three worksheets. The numbers were right: radius within $0.5\%$ of the video, rate unchanged.
Trajectory and centre moved \emph{together}, so every check agreed.}
\end{center}
\takeaway{No layer caught this. A frame of reference cannot be audited from inside itself.}
\end{frame}

% =========================================================================
\begin{frame}{What the layers caught --- and one they did not}
\vfill
\begin{center}
\small
\renewcommand{\arraystretch}{1.5}
\begin{tabular}{@{}p{4.2cm}p{6.6cm}@{}}
\toprule
\rowcolor{cInfra!12}
\textbf{Layer} & \textbf{What it caught} \\
\midrule
deterministic gate & ungrounded rate-vs-period wording, $4$ of $15$ \\
LLM judge (rater B) & \textbf{turntable-2: $\omega^2 r$ printed as $4.91$ when it is $4.31$} \\
phone sensor & the three teleports, $25.5\%$ on peak rate \\
a second measurement principle & the old fan method reporting $4.69$~rad/s at rest \\
\addlinespace[2pt]
\rowcolor{cAccent!10}
\textbf{none of them} & \textbf{the overlay $328$~px off its object --- found by opening the video} \\
\bottomrule
\end{tabular}
\end{center}
\vfill
\begin{center}
{\footnotesize The gate \emph{passed} the worksheet the judge caught --- $4.91$ does trace to
the measurement. It is the mean of $\omega^2 r$, not $\omega^2 r$ of the mean.}
\end{center}
\takeaway{Each layer sees a question the others cannot ask --- and none of them asks whether the
picture matches the world.}
\end{frame}

% =========================================================================
\begin{frame}{What is still unproven}
\vfill
\begin{center}
\small
\renewcommand{\arraystretch}{1.5}
\begin{tabular}{@{}p{5.6cm}p{5.2cm}@{}}
\toprule
\rowcolor{cAccent!10}
\textbf{Open} & \textbf{What would settle it} \\
\midrule
\textbf{Does anybody learn from it?} & the teacher panel \\
No clip has a sensor check any more & re-shoot one clip logging its gyroscope \\
Grounding asks two questions & split the criterion, then re-rate \\
Metre scale rests on one clip & a known length in every scene \\
Nothing rejects a teleport & build the guard; the sweep is still by hand \\
\bottomrule
\end{tabular}
\end{center}
\vfill
\begin{center}
{\footnotesize Withdrawing turntable-3 removed the only clip with a paired sensor log. That is
a real cost of the decision, not a footnote.}
\end{center}
\takeaway{Everything measured so far is a property of the text and the measurement. The first
row is the only one that answers the question the project exists to ask.}
\end{frame}

\end{document}
